%==============================================================================
%  TRINITY -- force-budget assessment, typeset as a standalone chapter
%
%  Build:  pdflatex trinity_pressure_chapter.tex   (x2, for the ToC)
%      or: latexmk -pdf trinity_pressure_chapter.tex
%
%  Equation / figure references use the source papers' own LaTeX labels
%  (grep the .tex files), not printed numbers.
%==============================================================================
\documentclass[11pt,a4paper,oneside]{report}

%--- encoding & fonts ---------------------------------------------------------
\usepackage[T1]{fontenc}
\usepackage[utf8]{inputenc}
\usepackage{mathpazo}                 % Palatino text + math
\usepackage[scaled=0.90]{helvet}      % sans for headings
\usepackage{courier}                  % typewriter
\usepackage{microtype}

%--- maths --------------------------------------------------------------------
\usepackage{amsmath,amssymb,bm}

%--- graphics -----------------------------------------------------------------
\usepackage{graphicx}

%--- layout -------------------------------------------------------------------
\usepackage[a4paper,left=3.1cm,right=3.1cm,top=2.7cm,bottom=2.9cm]{geometry}
\usepackage{fancyhdr}
\usepackage{titlesec}
\usepackage{enumitem}
\usepackage{booktabs}
\usepackage{array}
\usepackage{caption}
\usepackage{needspace}
\usepackage[dvipsnames,table]{xcolor}
\usepackage[most]{tcolorbox}
\usepackage{listings}
\usepackage[hidelinks]{hyperref}

%--- colours ------------------------------------------------------------------
\definecolor{tqink}{HTML}{1B2430}
\definecolor{tqaccent}{HTML}{9C2A2A}   % TRINITY red
\definecolor{tqblue}{HTML}{1F4E79}     % literature blue
\definecolor{tqrule}{HTML}{C8CDD4}
\definecolor{tqsoft}{HTML}{F4F5F7}
\definecolor{tqgreen}{HTML}{2E6B4F}

\hypersetup{colorlinks=true,linkcolor=tqblue,urlcolor=tqblue,citecolor=tqblue,
            pdftitle={TRINITY's force budget in context},
            pdfauthor={Assessment for Jia Wei Teh}}

%--- headings -----------------------------------------------------------------
\titleformat{\chapter}[display]
  {\normalfont\sffamily\bfseries\color{tqink}}
  {\raggedright\normalsize\sffamily\bfseries\color{tqaccent}\MakeUppercase{\chaptertitlename\ \thechapter}}
  {6pt}{\Huge\raggedright}[\vspace{4pt}{\color{tqrule}\titlerule[1.1pt]}]
\titlespacing*{\chapter}{0pt}{-20pt}{28pt}

\titleformat{\section}{\normalfont\sffamily\Large\bfseries\color{tqink}}{\thesection}{0.7em}{}
\titleformat{\subsection}{\normalfont\sffamily\large\bfseries\color{tqink}}{\thesubsection}{0.7em}{}
\titleformat{\subsubsection}{\normalfont\sffamily\normalsize\bfseries\color{tqink}}{}{0em}{}
\titlespacing*{\section}{0pt}{16pt}{6pt}
\titlespacing*{\subsection}{0pt}{12pt}{4pt}

\pagestyle{fancy}
\fancyhf{}
\renewcommand{\headrulewidth}{0.4pt}
\fancyhead[L]{\footnotesize\sffamily\color{tqink}TRINITY force budget vs.\ the literature}
\fancyhead[R]{\footnotesize\sffamily\thepage}

\setlength{\parskip}{0.45\baselineskip}
\setlength{\parindent}{0pt}
\setlist[itemize]{leftmargin=1.25em,itemsep=2pt,topsep=3pt}
\setlist[enumerate]{leftmargin=1.55em,itemsep=3pt,topsep=3pt}

\captionsetup{font=small,labelfont={sf,bf},labelsep=period}
\renewcommand{\arraystretch}{1.18}

%--- listings -----------------------------------------------------------------
\lstdefinestyle{tqpy}{
  language=Python, basicstyle=\ttfamily\footnotesize,
  keywordstyle=\color{tqblue}\bfseries, commentstyle=\color{tqgreen}\itshape,
  stringstyle=\color{tqaccent}, showstringspaces=false,
  backgroundcolor=\color{tqsoft}, frame=leftline, framerule=1.6pt,
  rulecolor=\color{tqrule}, xleftmargin=10pt, framexleftmargin=10pt,
  breaklines=true, columns=fullflexible, aboveskip=8pt, belowskip=6pt}
\lstset{style=tqpy}

%--- boxes --------------------------------------------------------------------
\newtcolorbox{keybox}[1][]{enhanced, breakable, colback=tqsoft, colframe=tqaccent,
  boxrule=0pt, leftrule=2.4pt, arc=1pt, left=9pt, right=9pt, top=7pt, bottom=7pt,
  fonttitle=\sffamily\bfseries\small, coltitle=white, #1}

\newtcolorbox{litquote}[1][]{enhanced, breakable, colback=white, colframe=tqblue,
  boxrule=0pt, leftrule=2.0pt, arc=0pt, left=9pt, right=9pt, top=6pt, bottom=6pt, #1}

\newtcolorbox{versionbox}[1][]{enhanced, breakable, colback=white, colframe=tqrule,
  boxrule=0.7pt, arc=1pt, left=9pt, right=9pt, top=7pt, bottom=7pt, #1}

%--- semantic macros ----------------------------------------------------------
% \_ inside \code/\lbl becomes a zero-width break opportunity, so long
% identifiers wrap instead of running into the margin.
\newcommand{\brkus}{\textunderscore\allowbreak}
\newcommand{\code}[1]{\texttt{\small\let\_\brkus #1}}
\newcommand{\lbl}[1]{\texttt{\footnotesize\let\_\brkus #1}}
\newcommand{\codet}[1]{\texttt{\scriptsize\let\_\brkus #1}}   % table-sized
\setlength{\emergencystretch}{3em}
\newcommand{\PI}{Paper~I}
\newcommand{\PII}{Paper~II}

\newcommand{\Rw}{\mathcal{R}_{w}}
\newcommand{\Ri}{\mathcal{R}_{i}}
\newcommand{\Rf}{\mathcal{R}_{f}}
\newcommand{\Rch}{R_{\mathrm{ch}}}
\newcommand{\Req}{R_{\mathrm{eq}}}
\newcommand{\RSt}{R_{\mathrm{St}}}
\newcommand{\pdotw}{\dot p}
\newcommand{\ap}{\alpha_p}
\newcommand{\Pb}{P_{b}}
\newcommand{\Pram}{P_{\mathrm{ram}}}
\newcommand{\PHII}{P_{\mathrm{HII}}}
\newcommand{\Pdrive}{P_{\mathrm{drive}}}
\newcommand{\Pconf}{P_{\mathrm{conf}}}
\newcommand{\Pext}{P_{\mathrm{ext}}}
\newcommand{\PCa}{P_{\mathrm{C3a}}}
\newcommand{\ci}{c_i}
\newcommand{\rhobar}{\bar\rho}
\newcommand{\Vw}{\mathcal{V}_w}
\newcommand{\Msh}{M_{\mathrm{sh}}}
\newcommand{\Fsp}{F_{b,\mathrm{Sp}}}
\newcommand{\Fcem}{F_{\mathrm{CEM}}}

\newcommand{\up}[1]{\textcolor{tqaccent}{\textbf{#1}}}
\newcommand{\dn}[1]{\textcolor{tqblue}{\textbf{#1}}}

%==============================================================================
%  BODY -- revision 2
%==============================================================================
\begin{document}
\setcounter{chapter}{0}

%------------------------------------------------------------------ front matter
\thispagestyle{empty}
\begin{center}
  \vspace*{2.0cm}
  {\sffamily\bfseries\Huge\color{tqink} The TRINITY force budget\\[4pt] in context\par}
  \vspace{10pt}
  {\color{tqrule}\rule{0.55\textwidth}{1.1pt}}\par
  \vspace{12pt}
  {\sffamily\large\color{tqaccent}
   $P_b$ \quad$\cdot$\quad $P_{\mathrm{ram}}$ \quad$\cdot$\quad
   $P_{\mathrm{HII}}$ \quad$\cdot$\quad $P_{\mathrm{drive}}$\par}
  \vspace{22pt}
  {\sffamily\normalsize\color{tqink}
   An assessment against Lancaster et al.\ \PI\ \& \PII,\\
   Geen et al.\ (2019), Geen \& de Koter (2022), and Haid et al.\ (2018)\par}
  \vspace{18pt}
  {\sffamily\normalsize\bfseries\color{tqaccent} Revision 2\par}
  \vspace{26pt}
  {\sffamily\small\color{tqink}
   Prepared for Jia Wei Teh \\[3pt]
   \today\par}
\end{center}

\vfill
\begin{versionbox}
\small
\textbf{\sffamily Revision 2 supersedes revision 1 of the same date.}
Revision 1 contained errors of substance, not only of detail.
Section~\ref{sec:errata} lists every one of them and what replaced it.
The four largest: \up{Section 4.1's $\ap$ diagnosis was wrong} --- TRINITY already
carries $\ap$, and $\ap \equiv 1$ in the momentum phase is the \emph{correct} 1D value,
not a defect; \up{the ``Haid--Lancaster tension'' does not exist}; \up{the Geen 2022 /
\textsc{warpfield} chronology was inverted}; and \up{the shell-mass recommendation was
contradicted by your own measurement} --- the momentum shell is ${\sim}99.5\%$ ionised,
not a skin.

\medskip
\textbf{\sffamily Provenance.}
Code state read: \code{\textasciitilde/unsync/Code/Trinity}, read-only --- nothing under
\code{trinity/} was modified. Papers read in full:
\code{lancaster2025.tex}, \code{lancaster2025b.tex}, \code{geen2019.tex},
\code{geen2022.tex}, \code{Haid118\_arxiv.tex}.
Workstream documents read at the 2026-08-18 pull:
\code{docs/dev/phii-identity/\{README.md, PLAN.md, LITERATURE\_ASSESSMENT.md\}},
including Batches~11 and 12 and the B11 cross-check.
Equation and figure references use the source papers' own \LaTeX{} labels, not printed
numbers, so every claim is checkable at source.
Section~\ref{sec:verify} records what was verified and how, what the second adversarial pass then
changed in revision 2's own draft (Section~\ref{sec:pass2}), and what could \emph{not} be verified.
\end{versionbox}

\clearpage
{\hypersetup{linkcolor=tqink}\tableofcontents}
\clearpage

%==============================================================================
\chapter{The TRINITY force budget in context}
%==============================================================================

\section{Verdict}

\begin{keybox}
\textbf{\sffamily The force budget is structurally sound and, on several specific points,
ahead of all five papers.} The C3c regime switch is defensible, and in the momentum phase
its crossover point is \emph{exactly} Lancaster's characteristic radius $\Rch$
\up{evaluated at $\ap=1$}, which is the value TRINITY's momentum phase carries --- so you reached
from a confinement argument the scale they reach from force balance.

\medskip
\textbf{\sffamily The two highest-value improvements are the two your own Batches 11 and 12
have already measured}, not the ones revision 1 nominated.
\begin{enumerate}[leftmargin=1.6em]
\item \textbf{The recombination balance volume.} \code{get\_phii\_c3c} balances
      $Q_{i,\mathrm{abs}}$ over $\tfrac{4\pi}{3}R_2^3$ --- the wind cavity.
      Lancaster \lbl{eq:ionreceq2}, Geen 2019 \lbl{wind:photoequilibrium}, Geen 2022
      \lbl{eqn:photoionisation\_equilibrium\_uniform} \emph{and TRINITY's own}
      \code{shell\_structure.py:243} all balance over $r_i^3-r_w^3$, i.e.\ over the volume
      C3a uses, \emph{excluded}. This is an inconsistency internal to TRINITY, not merely a
      difference from the literature, and it is what your seam-C mass double-book measures.
\item \textbf{Shell inertia.} Subtracting the ionised mass is worth $+8.55\%$ / $+9.22\%$
      in $R_2(t{=}1.5)$ at nominal wind (B11.C2, against a control good to $0.871\%$).
      Lancaster \PI\ already adopts this form (\lbl{eq:pr\_spitzer\_adj}), with its own
      consistency caveat.
\end{enumerate}

\medskip
\textbf{\sffamily $\ap$ is not the answer to the momentum-phase question}, and revision 1
was wrong to say it was. Section~\ref{sec:alphap} replaces that recommendation with the
correct reading of what TRINITY already computes.
\end{keybox}

%==============================================================================
\section{What TRINITY actually computes}
\label{sec:what}

Verified against the staged source; every file:line was re-checked this revision.

\begin{table}[htbp]
\centering
\small
\begin{tabular}{@{}>{\raggedright\arraybackslash}p{0.20\textwidth}
                  >{\raggedright\arraybackslash}p{0.31\textwidth}
                  >{\raggedright\arraybackslash}p{0.40\textwidth}@{}}
\toprule
\textbf{\sffamily Quantity} & \textbf{\sffamily Source} & \textbf{\sffamily Expression}\\
\midrule
$\Pb$ (1a/1b) & \codet{get\_bubbleParams.bubble\_E2P} &
  $(\gamma-1)E_b/[\tfrac{4\pi}{3}(R_2^3-R_1^3)]$\\
$R_1$ & \codet{get\_r1} / \codet{solve\_R1} &
  root of $\sqrt{L_{\mathrm{mech}}(R_2^3-R_1^3)/(v_{\mathrm{mech}}E_b)}=R_1$\\
$\Pb$ \emph{as the ODE sees it} & \codet{get\_effective\_bubble\_pressure}, \codet{:495--503} &
  same, but $R_1\!\to\!\frac{t-t_{\mathrm{SF}}}{10^{-3}}R_1$ for $t\le t_{\mathrm{SF}}+10^{-3}$\,Myr\\
$\Pram$ & \codet{get\_bubbleParams.pRam} &
  $L_{\mathrm{mech}}/(2\pi R_2^2 v_{\mathrm{mech}})$, $v_{\mathrm{mech}}=2L_{\mathrm{mech}}/\pdotw_{\mathrm{tot}}$\\
$\PHII$ & \codet{get\_phii\_c3c} &
  $\frac{\mu_c}{\mu_i}k_BT_i\sqrt{3Q_{i,\mathrm{abs}}/(4\pi\chi_e\alpha_B R_2^3)}$ if $>\Pconf$,
  else \up{exactly 0} (\codet{:365})\\
$\Pconf$ & \codet{get\_bubbleParams.py:365} &
  \codet{params['Pb']} --- bubble \emph{thermal} pressure in 1a/1b/1c, the \emph{ram}
  pressure in phase 2 (\codet{run\_momentum\_phase.py:585,669,891})\\
$\Pext$ & \codet{energy\_phase\_ODEs.get\_press\_ion} &
  $\frac{\mu_c}{\mu_i}n(r_{\mathrm{sh}})k_BT_i$ when $f_{\mathrm{abs,ion}}<1$;
  $+\,n_{\mathrm{ISM}}k_BT$ beyond $r_{\mathrm{cloud}}$\\
$\Pdrive$ (1a) & \codet{energy\_phase\_ODEs.py:388} & $\max(\Pb,\PHII)$\\
$\Pdrive$ (1b) & \codet{run\_energy\_implicit\_phase.py:532} & $\max(\Pb,\PHII)$\\
$\Pdrive$ (1c) & \codet{run\_transition\_phase.py:331} & $\max(\Pb,\PHII+\Pram)$\\
$\Pdrive$ (2)  & \codet{run\_momentum\_phase.py:265,445} & $\PHII+\Pram$\\
Shell momentum & \codet{energy\_phase\_ODEs.py:263} &
  $\Msh\dot v_2 = 4\pi R_2^2(\Pdrive-\Pext)-\dot\Msh v_2-F_{\mathrm{grav}}+F_{\mathrm{rad}}$\\
\bottomrule
\end{tabular}
\caption{The pressure terms and where they are set.}
\end{table}

\begin{keybox}
\textbf{\sffamily Which of these the solver actually integrates.} This matters for any patch scoped
by the table. The transition phase's RHS is \code{get\_ODE\_transition\_pure}
(\code{run\_transition\_phase.py:630}), which delegates the momentum equation to
\code{get\_ODE\_Edot\_pure} (\code{:231-233}). So the \up{live} drive expressions are
\code{energy\_phase\_ODEs.py:253} (the \code{current\_phase == 'transition'} branch) and \code{:256}
(the \code{else} branch), plus \code{run\_momentum\_phase.py:445}. \up{Three live sites.}
\code{run\_transition\_phase.py:331}, \code{run\_energy\_implicit\_phase.py:532},
\code{energy\_phase\_ODEs.py:388} and \code{run\_momentum\_phase.py:265} all sit in
\code{compute\_forces\_pure} / \code{compute\_derived\_quantities} and are \dn{reporting only} ---
they set \code{params['P\_drive']} for the snapshot and never reach the integrator.
\code{energy\_phase\_ODEs.py:385} is genuinely unreachable: \code{compute\_derived\_quantities} is
called from one place (\code{run\_energy\_phase.py:242}) with \code{current\_phase == 'energy'}, so
its \code{'transition'} branch never fires. The formulas are identical either way, so the table is
right as physics; it is the \emph{edit surface} that differs.
\end{keybox}

Two identities follow, both verified symbolically (Section~\ref{sec:verify}).

\subsection*{Identity I --- the $R_1$ closure is free-wind ram-pressure balance, so the
code already carries a momentum enhancement factor}

Substituting the \code{get\_r1} root into \code{bubble\_E2P},
\[
\Pb \;=\; (\gamma-1)\,\frac{3}{4\pi}\,\frac{L_{\mathrm{mech}}}{v_{\mathrm{mech}}R_1^2}
\;\;\xrightarrow{\;\gamma=5/3\;}\;\;
\frac{L_{\mathrm{mech}}}{2\pi v_{\mathrm{mech}}R_1^2}\;=\;\frac{\pdotw}{4\pi R_1^2},
\]
so the wind force reaching the shell is
$F = 4\pi R_2^2 \Pb = \pdotw\,(R_2/R_1)^2$.

\medskip
\emph{Two caveats.} The clean $\pdotw/4\pi R_1^2$ form is exact only at the default
\code{gamma\_adia = 5/3}, which is user-exposed. And the identity holds for the
\up{stored} \code{params['Pb']} at all times, but for the pressure the \up{ODE actually
integrates} only \emph{outside} the \code{dt\_switchon} ramp window: for
$t \le t_{\mathrm{SF}}+10^{-3}$\,Myr the ODE uses $R_1$ scaled linearly toward zero
(\code{get\_bubbleParams.py:495--503}). With \code{TFINAL\_ENERGY\_PHASE = 3e-3} that is up
to the first \emph{third} of phase 1a's maximum duration (exactly a third at the default
\code{tSF = 0}; the phase-1a loop bound is absolute $t$ while the ramp is $t_{\mathrm{SF}}$-relative),
and your own \code{PLAN.md} \S1(3) measures the two pressures differing by up to $3.31\times$
inside it. Any $\ap$ diagnostic
built on Identity I must be reported as \emph{invalid} there rather than silently wrong.

\subsection*{Identity II --- the momentum phase applies exactly $F=\pdotw$, i.e.\ $\ap=1$ in the
force sense}

This is \emph{not} a consequence of Identity I. In phase 2 \code{bubble\_E2P} is never called; the
drive is built directly from \code{pRam(R2)} (\code{run\_momentum\_phase.py:445,272,585,669}), and
with $v_{\mathrm{mech}}=2L_{\mathrm{mech}}/\pdotw_{\mathrm{tot}}$ (\code{update\_feedback.py:181}),
$4\pi R_2^2 \Pram = \pdotw_{\mathrm{tot}}$ identically.

\code{run\_momentum\_phase.py:587--588} (and \code{:893}) additionally set
\code{params['R1']} to \code{R2}, with the comment \emph{``Set R1 = R2 (no inner shock in momentum
phase)''}. \up{That assignment is bookkeeping, not dynamics} --- \code{params['R1']} is read in
exactly one place in the tree --- \code{get\_bubbleParams.py:115}, inside
\code{cool\_beta\_to\_Ebdot}, which is energy/implicit-phase code. It sets the output column and states the modelling assumption;
it does not force anything.

The assumption it states is nonetheless the right one, and Lancaster say the same from the other
side:

\begin{litquote}
\small ``\ldots in the limit that $\ap \to 1$ we have $\Rf \to \Rw$ and thus that, in the
idealized, purely momentum-driven solution the entire bubble is made up of the free wind
region.''\hfill{\scriptsize\PI, \lbl{subsubsec:ec\_wind}}
\end{litquote}

So TRINITY's momentum phase \emph{is} the $\ap = 1$ idealised momentum-driven limit ---
not a missing factor. See Section~\ref{sec:alphap}.

%==============================================================================
\section{Where you sit relative to each paper}

\subsection{Lancaster \PI\ --- the direct comparison}
\label{sec:lanc1}

\PI\ is analytic; it runs no simulations of its own (the 3D runs are \PII's). It criticises
\textsc{warpfield}-class models in at least nine passages, which reduce to \up{six} distinct
criticisms --- not the three revision 1 listed.

\begin{table}[htbp]
\centering\footnotesize
\begin{tabular}{@{}c>{\raggedright\arraybackslash}p{0.60\textwidth}
                  >{\raggedright\arraybackslash}p{0.24\textwidth}@{}}
\toprule
\textbf{\sffamily \#} & \textbf{\sffamily Criticism (labels)} & \textbf{\sffamily TRINITY}\\
\midrule
1 & \textbf{The thin-shell lumping itself.} Channels are ``combine[d] many feedback channels
    artificially into a single thin-shell evolution equation\ldots This approximation is made in
    order to simplify calculations but does not have rigorous theoretical foundations''
    (\lbl{sec:intro}); ``The model assumes that the bubble is made up of a single thin-shell where
    all forces\ldots are applied'' (\lbl{subsec:theory\_review}); it ``prevent[s] a faithful
    representation of the gas density and ionization structure'' (\lbl{subsec:problems})
  & \dn{open} --- \PI's central thesis; \S\ref{sec:volume}, \S\ref{sec:closure} answer it. Partial
    credit on the third: the shell solver \emph{does} resolve the structure, but does not separate
    $\Rw$ from $\Ri$\\
2 & \textbf{No turbulently enhanced cooling.} ``does not include a model for turbulently enhanced
    cooling\ldots which could make the WBBs act in a more momentum-driven manner much earlier''
    (\lbl{subsec:theory\_review}); conduction-only, spherical dissipation misses the dominant loss
    channel (\lbl{subsec:problems})
  & \dn{open}; what the \code{cooling\_boost} knobs attack --- see below\\
3 & \textbf{The photoionised thermal-pressure force is omitted from the momentum equation}
    (\lbl{subsec:theory\_review}) & \up{structurally fixed}\\
4 & \textbf{Constant-density PIR.} ``these semi-analytic models must take into account
    inhomogeneities in the background medium into which they expand'' (\lbl{subsec:problems});
    ``Principal among these would be to remove the assumption of a constant density PIR''
    (\lbl{sec:conclusion}) & \dn{open} --- \S\ref{sec:thin}'s $f_{\mathrm{ion}}$ knob is the cheap
    1D proxy\\
5 & \textbf{The early independent phase} is ``distinct from simpler models\ldots which use a single
    thin-shell approximation'' (\lbl{subec:early\_evol}, the paper's own typo) & \dn{open} --- needs
    $\Rw\ne\Ri$. Note the paper says ``distinct from'', not ``cannot represent''\\
6 & \textbf{LyC ``trapping'' is assumed.} It is ``convenient for theoretical modeling\ldots [but]
    is often not seen in simulations with a cluster of massive stars\ldots or uniform density
    backgrounds'' (\lbl{app:trapping}) & partly answered --- you compute trapping rather than
    assuming it; the geometry that produces it is still 1D\\
\bottomrule
\end{tabular}
\caption{\PI's criticisms of \textsc{warpfield}, grouped and scored against TRINITY. Revision 1
listed three of these, omitting \PI's central thesis. A seventh criticism in the same section ---
that instantaneous thermal/ram-pressure balance ``ignores the inertia of the shell''
(\lbl{subsec:theory\_review}, citing Raga et al.\ 2012b) --- is aimed at the pressure-balance class,
which TRINITY is not in; it appears in Section~\ref{sec:geen19} against Geen 2019 instead.}
\end{table}

\subsubsection*{On criticism 4}
\code{include\_PHII} + \code{get\_phii\_c3c} put a photoionised term into $\Pdrive$ in all
four phases, and in the transition and momentum phases it dominates in every configuration
measured. In 1a/1b the confined branch fires and $\PHII$ is exactly \code{0.0}, so
$\Pdrive=\Pb$ there --- \up{but that is a property of the regime, not a theorem.} Batch 7's
confinement screen records \code{B3MW001} ($L_w\times0.01$) at \up{78.4\% HII-dominated in
the energy phase}, with \code{ratio\_max} = 4.927 (\code{data/b7\_confinement\_screen.csv}; that run
is \code{run\_complete = False} and never reaches transition or momentum, so \code{PLAN.md} marks it
VOID for any \emph{driving-branch} claim --- it is evidence about 1a/1b only). Revision 1's ``0.0000 on every energy and implicit row of all five configs''
was a stale Batch-5 read and is withdrawn.

\subsubsection*{On criticism 6}
\PI's ``Prospects'' sentence is easy to truncate favourably. In full:

\begin{litquote}
\small ``The ideal version of these semi-analytic models would include a parameterized model
for this heat dissipation that is solved on-the-fly and included in the WBB energy
equation\ldots This is done for the case of only conductive heat dissipation in a spherical
scenario by the \texttt{WARPFIELD} models, \textbf{but as Lancaster et al.\ (2021a,c) have shown, cooling in
turbulently-mixed intermediate-temperature gas can certainly dominate energy losses.}''\hfill{\scriptsize\PI, \lbl{subsec:problems}}
\end{litquote}

So TRINITY has the \emph{architecture} \PI\ asks for and half the \emph{physics}. That
framing is defensible; the sentence must be quoted past the comma.

\subsubsection*{The C3c switch point is Lancaster's $\Rch$ --- exactly, and only in the
momentum phase}

Setting $\PCa(R_2)=\Pconf(R_2)$ with $\Pconf=\Pram$:
\[
\frac{\pdotw}{4\pi R_2^2}=\rhobar \ci^2\!\left(\frac{\RSt}{R_2}\right)^{3/2}
\;\;\Longleftrightarrow\;\;
R_2=\frac{\Req^4}{\RSt^3}\;=\;\Rch ,
\]
using $\Req^2 \equiv \ap\pdotw/(4\pi\rhobar \ci^2)$ (\up{\lbl{eq:RE\_MD\_def}}, restated in
\lbl{eq:ReqMD\_app}; note \lbl{eq:eta\_def} defines $\zeta$ itself, not $\Req$) and the relation
$\Rch=\Req^4/\RSt^3$ (\up{\lbl{eq:Rch\_Req\_rel}}, restated dimensionally in
\lbl{eq:Rch\_app}; the \emph{definition} of $\Rch$ is \lbl{eq:Rch\_def},
$\Rch\equiv\frac{\alpha_B}{12\pi(\mu_Hm_H\ci^2)^2}\frac{\ap^2\pdotw^2}{Q_0}$). SymPy returns
crossover$/\Rch = 1$ exactly \up{at $\ap=1$}; since $\Rch\propto\ap^2$, at any other $\ap$ the
crossover sits at $\Rch/\ap^2$. That is a consistency check rather than a caveat: TRINITY's phase-2
confining pressure $\pdotw/(4\pi R_2^2)$ \emph{is} the $\ap=1$ wind pressure, so $\ap=1$ is the
right value to evaluate $\Rch$ at. So in phase 2 your confined branch is precisely Lancaster's
wind-dominated $R_2<\Rch$ and your unconfined branch their PIR-dominated $R_2>\Rch$.

\medskip
\textbf{Three caveats.}
\begin{itemize}
\item \code{get\_phii\_c3c} compares against \code{params['Pb']}, which is $\Pram$
      \up{only in the momentum phase}. In 1a/1b/1c it is the bubble \emph{thermal} pressure
      --- larger by $(R_2/R_1)^2$ --- so there the switch is the energy-driven analogue of
      $\Rch$: the same force-balance \emph{structure}, but not the quantity Lancaster
      tabulate (they define $\Rch$ for the MD case). A natural generalisation, not an
      error --- but \code{t\_cross} in your ladder is \up{not} the $R_2=\Rch$ crossing.
\item The reduction assumes $f_{\mathrm{abs,ion}}=1$. Since $\PCa\propto\sqrt{f_{\mathrm{abs}}}$,
      $\PCa\propto R_2^{-3/2}$ and $\Pram\propto R_2^{-2}$, the crossover scales as
      $R_2^{\mathrm{cross}}\propto f_{\mathrm{abs}}^{-1}$: \up{partial absorption pushes the
      crossover outward}, it does not shrink it. (Revision 1 said $f_{\mathrm{abs}}^{1/3}$,
      wrong in both exponent and sign.) In practice this rarely bites --- B11.0 measures
      $f_{\mathrm{abs}}=1.0000$ on 16/16 transition and 13/17 momentum driving rows at
      nominal wind, and on \emph{all} driving rows at $L_w\times0.1$.
\item $\Rch \propto \ap^2$, so quoting a $R_2/\Rch$ requires stating which $\ap$ was used.
\end{itemize}

With those caveats the correspondence stands and is worth stating in the method paper.
Batch 8's result --- C3a reproduces Hosokawa--Inutsuka to 0.0000\% over $R/\RSt\in[2,50]$ --- is
the same statement seen from the Spitzer side, and should be quoted with the caveat \code{PLAN.md}
attaches to it: \emph{``Not independent confirmation --- HI is derived from the same momentum
equation, so once the algebra gates hold the ODE must return HI.''} (Its G8.4 gate failed as
registered, 9.511\% against a 5\% bar, and was amended.)

\medskip
\textbf{Where TRINITY is better than the CEM:} the CEM ignores gravity, external pressure,
direct and indirect radiation pressure, the cloud density profile, and time-variable
feedback. It \emph{does} solve an energy equation in its ED variant
(\lbl{subsec:ed\_jfb}), so the honest count is \up{six}, not seven, and what the ED-CEM
lacks is a self-consistent \emph{dissipation model}, not the equation.

\textbf{Where the CEM is better:} it separates $\Rw$ from $\Ri$. TRINITY's $R_2$ does both
jobs. That is the root of \S\ref{sec:volume} and \S\ref{sec:closure}.

%------------------------------------------------------------------------------
\subsection{Lancaster \PII\ --- the calibration, and what it is a calibration \emph{of}}
\label{sec:lanc2}

\PII's \lbl{tab:cem\_comp} reports $\langle\ap\rangle$ = 4.64 / 4.78 / 4.57 (\code{HWR} at
$N=128/256/512$) and 5.57 / 6.20 / 6.82 (\code{MWR}). \up{The MWR values rise monotonically with
resolution, and the paper reports the trend without claiming convergence}: \emph{``For the}
\code{WRM} \emph{simulations, we see a tendency to larger $\ap$ with higher resolution.''} (It also
notes \emph{``there is no clear pattern towards smaller values of $\langle\Delta\rangle$ with
resolution.''}) Quote the range, not a single number.

Crucially, these $\ap$ are measured from \lbl{eq:alphap\_derive},
$\ap = \tfrac34\frac{\Vw/4}{\langle v_{\mathrm{out}}\rangle}\frac{4\pi\Rw^2}{A_w}$, which is a
product of \up{two} independent factors, and \PII\ names them as such: \emph{``the dynamics of WBBs
are affected by properties of cooling at their interfaces in two distinct ways: dissipation and
geometry.''}

\begin{itemize}
\item the \textbf{dissipation} factor $\tfrac{3}{16}\Vw/\langle v_{\mathrm{out}}\rangle$ --- how much
      energy the interface loses;
\item the \textbf{geometry} factor $4\pi\Rw^2/A_w$ --- the spherical-to-fractal area ratio.
\end{itemize}

\up{Only the second is unavailable in 1D}, where $A_w\equiv4\pi\Rw^2$ identically. The first is a
physics parameter a 1D code can and does carry --- TRINITY's is the bubble energy budget, and it is
why $(R_2/R_1)^2\gg1$ in the energy phase. So the correct statement is \emph{not} ``1D forbids
$\ap>1$'' (revision 1's error, in the opposite direction from its first one) but: \up{1D forbids the
geometric half of the enhancement, and $A_w$ is exactly what \code{cooling\_boost\_fA}
parameterises on the energy side}. \PII\ is in any case explicit that the values are resolution- and
geometry-dependent: \emph{``The exact values of $A_w$ and $\langle v_{\mathrm{out}}\rangle$ (and
therefore $\ap$) in reality will likely be different from those in our simulation.''}

\begin{keybox}
\textbf{\sffamily $(R_2/R_1)^2$ \emph{is} the directly comparable quantity --- revision 1
said the opposite and was wrong.}

Lancaster take the hot-gas pressure at the post-shock value
$P_{\mathrm{hot}} = 3\pdotw/(16\pi\Rf^2)$, while TRINITY's \code{get\_r1} uses the free-wind
\emph{ram} value $\pdotw/(4\pi R_1^2)$. At the same pressure $\Rf = (\sqrt3/2)R_1$, so
$x \equiv \Rw/\Rf = (2/\sqrt3)(R_2/R_1)$, and substituting into \lbl{eq:alphap\_shock},
$\ap=\tfrac14[3x^2+x^{-2}]$:
\[
\ap \;=\; \left(\frac{R_2}{R_1}\right)^{2} + \frac{3}{16}\left(\frac{R_1}{R_2}\right)^{2}.
\]
The $4/3$ convention mismatch cancels the $3/4$ in the leading term \emph{exactly}. The residual
$\tfrac{3}{16}(R_1/R_2)^2$ is $18.75\%$ of the leading term at $R_2/R_1=1$, $1.2\%$ at 2,
$0.2\%$ at 3, negligible above.
\end{keybox}

\textbf{Two things must be said about this mapping, or it misleads.} First, it comes from matching
the two conventions \emph{at the same bubble pressure} --- the right invariant, because the bubble
pressure is what exerts the force on the shell, and $R_1$ and $\Rf$ denote the same physical surface
(the wind termination shock) computed under two shock conditions: TRINITY's the ram-pressure form,
Lancaster's the exact strong-shock post-shock value. Identifying $\Rf\equiv R_1$ \emph{geometrically}
instead would give $\tfrac34(R_2/R_1)^2+\tfrac14(R_1/R_2)^2$ and no cancellation; the
pressure-matched version is the one to use.

Second, \up{$\ap$ and the force ratio are not the same quantity}. Lancaster's own
\lbl{eq:Phot\_EC} gives the exact force as $4\pi\Rw^2 P_{\mathrm{hot}} = \tfrac34\pdotw x^2$, and
writes $\approx\ap\pdotw$ only \emph{``using the assumption $\ap \gtrsim 1$''} --- i.e.\ dropping the
$x^{-2}$ term. Under the pressure match, $\tfrac34 x^2 = (R_2/R_1)^2$ exactly:
\[
\underbrace{(R_2/R_1)^2}_{\text{TRINITY's exact }F/\pdotw}
\;=\;\underbrace{\tfrac34 x^2}_{\text{Lancaster's exact }F/\pdotw}
\;\le\;\underbrace{\tfrac14[3x^2+x^{-2}]}_{\ap,\ \lbl{eq:alphap\_shock}} .
\]

\begin{keybox}
\up{Emit $(R_2/R_1)^2$.} It is TRINITY's exact force ratio \emph{and} Lancaster's exact force ratio,
and it is the quantity \PII's tabulated $\ap$ approximates. Adding the $\tfrac{3}{16}(R_1/R_2)^2$
term reproduces \lbl{eq:alphap\_shock} exactly but is actively misleading where $R_2\to R_1$: it
would report $1.1875$ for a bubble whose force ratio is exactly 1.
\end{keybox}

One further limit on the comparison: \lbl{eq:alphap\_shock} is flagged as true \emph{``excluding
geometric effects which are treated more carefully in Appendix A of [Lancaster 2024a]''}, so a 1D
$(R_2/R_1)^2$ maps onto the \up{spherical-equivalent} $\ap$ only. \PII\ itself performs exactly this
inversion for the literature in \lbl{app:sim\_comp}, so the procedure is theirs, not an invention.

\medskip
\textbf{Photoionisation raises $\ap$} --- 2.55 $\to$ 4.66 (HD), 4.09 $\to$ 6.20 (MHD) --- by
removing collisional Ly$\alpha$ as an interface coolant and by smoothing background
inhomogeneity so $A_w$ drops. Note the sign: adding photoionisation makes the \emph{wind}
more effective. One caveat worth carrying: \PII\ \lbl{subsec:cooling} documents that the
wind-only baseline suffers numerical diffusion of neutral H into $10^4$--$10^{5.5}$\,K gas,
\emph{``leading to an excess of cooling due to Ly$\alpha$ that otherwise would not be
present in this gas''}. The authors argue this \emph{``does not detract from any of the main results''}
and treat the increase as physical; they do \up{not} attribute any part of
$2.55 \to 4.66$ to the artefact. Cite the increase as their result, and the caveat as theirs
too.

%------------------------------------------------------------------------------
\subsection{Geen \& de Koter 2022 --- the same shell physics, with the chronology the right
way round}
\label{sec:geen22}

\code{get\_shellODE.py} implements Geen 2022's \lbl{eqn:draine1}--\lbl{eqn:draine3} --- the
Draine (2011) hydrostatic dusty-ionised-shell system with $\phi$ and $\tau$ integrated
outward --- with the same inner boundary condition
$P_w = P_i \equiv (m_H/X)n_i\ci^2$ (\lbl{eqn:PwPibalance}; TRINITY
\code{shell\_structure.py:125-126}). Geen 2019 closes the same system identically
(\lbl{wind:windpressurebalance}).

\begin{keybox}
\textbf{\sffamily This is a common inheritance, not TRINITY implementing Geen.}
Geen 2022 routes \emph{its own} equations through Draine (2011) and
Mart\'inez-Gonz\'alez et al.\ (2014); \PI\ attributes \textsc{warpfield}'s hydrostatic shell to
Abel+2005 / Pellegrini+2007 / Draine 2011 / Kim+2016, so the shared ancestor is Draine (2011), not
the Mart\'inez-Gonz\'alez line. Geen 2022 then cites \code{Rahner2017} --- \textsc{warpfield},
TRINITY's parent --- in the list of models that already apply this analysis with an embedded wind
bubble. Revision 1's \emph{``you already
implement their model''} inverted a chronology in which TRINITY's lineage is the
\emph{earlier} work. The correct statement is that the two agree term-for-term, which is a
mutual corroboration.
\end{keybox}

\textbf{What TRINITY does that they do not:} solve it self-consistently along a dynamical
trajectory. Geen 2022 solve it on a prescribed $r_w(t)$ from their analytic Weaver-like
solution (\lbl{eqn:rwt}).

\textbf{What they do that TRINITY does not} --- revision 1's ``you are strictly ahead here''
is withdrawn:
\begin{itemize}
\item \textbf{Metallicity-dependent dust:} $\sigma_d = 10^{-21}\,\mathrm{cm}^{-2}\,Z/Z_\odot$,
      explicitly.
\item \textbf{Star-dependent $T_i$:} tabulated per stellar model
      (\lbl{table:app\_starprops}), not pinned to $10^4$\,K.
\item \textbf{$\Omega < 4\pi$ solid-angle geometry.} (\code{coverFraction} is an energy
      leak, not a solid angle --- different physics.)
\item \textbf{An analytic overflow radius}, \lbl{eqn:overflowcondition}, with the clean
      $\omega$ vs $5/4$ threshold: for $\omega<5/4$ overflow becomes \emph{less} likely with
      radius (trapping wins), for $\omega>5/4$ \emph{more} likely. \code{densPL\_alpha} spans
      both sides, so this is a sharp, cheap validation target.
\end{itemize}
The first two are exactly what Section~\ref{sec:metal} charges TRINITY with lacking.

\medskip
\textbf{The overflow criteria are not the same criterion.} Geen's \emph{numerical} test is
$M_i \ge M(<r_i)$ with $M_i=\int_{r_w}^{r_i}\Omega r^2n_i(m_H/X)\,dr$
(\lbl{eqn:photoionised\_mass}); their \emph{analytic} version adds \emph{``a further
simplifying assumption that the mass swept up by the shell $M(<r_i)\simeq M(<r_w)$.''}
TRINITY's loop (\code{shell\_structure.py:158,182}) runs
\code{while not is\_allMassSwept and not is\_phiDepleted} against
\code{mShell\_end = params['shell\_mass']} (\code{:107}), integrating outward from
\code{rShell0 = R2} (\code{:106}). \code{shell\_mass} is set from \code{get\_mass\_profile(R2,
params)} immediately before each shell solve in every runner, i.e.\ it is the mass originally
interior to $R_2$ --- $M(<r_w)$ in Geen's notation. (B11.0's \code{shell\_mass/M\_avail} =
0.999997--1.000000 while $R_2$ grows $4.65\times$ is independent confirmation: only $M(<R_2)$ can
saturate that way.) \up{So TRINITY implements Geen's \emph{analytic}
criterion, not their numerical one}, compares against the smaller reference mass, and
therefore overflows \emph{earlier}. The discrepancy is $M(<r_i)/M(<r_w) = (r_i/r_w)^{3-\omega}$;
Geen 2022's \lbl{fig:ionisedshellthickness} shows $(r_i-r_w)/r_w$ approaching unity near
overflow, so at $\omega=2$ this is a factor up to ${\sim}2$ in mass. Worth a one-line
comparison in the paper --- not obviously wrong, but it is a choice you currently make
implicitly.

%------------------------------------------------------------------------------
\subsection{Geen et al.\ 2019 --- you are better on the one thing that matters}
\label{sec:geen19}

Geen 2019's dynamical closure is a \emph{ram-pressure balance} at the front,
$n_i\ci^2 = n(r_i)(\dot r_i+v_0)^2$ (\lbl{photo:externalpressure}), with \up{no shell
inertia}. Lancaster \PI\ \lbl{subsec:theory\_review} flags instantaneous force balance as
the key weakness of this class. \up{TRINITY integrates the full thin-shell momentum equation
including $\dot\Msh v_2$ and $\Msh\dot v_2$} --- Hosokawa--Inutsuka rather than Spitzer.
That is the right call and you should say so.

Two things to take from Geen 2019:
\begin{itemize}
\item \textbf{It licenses your transparent cavity.}
      \begin{litquote}
      \small ``Since the temperature of this gas is typically well above the limit to be
      collisionally ionised, the UV photons from the star are not absorbed by the wind
      bubble.''\hfill{\scriptsize Geen 2019, \lbl{winds\_in\_uv}}
      \end{litquote}
      TRINITY's \code{phi0 = 1} at \code{shell\_structure.py:119-120} is therefore the
      standard picture and is \up{correct}. What is \emph{not} standard is the balance
      \emph{volume} (Section~\ref{sec:volume}).
\item \textbf{A cross-check quantity.}
      $C_w = 2^{1/4}\big(\frac{\pdotw_w}{4\pi}\frac{X}{m_H}\frac{1}{\ci^2}\big)^{3/2}
      \big(\frac{Q_H}{\alpha_B}\frac{3}{4\pi}r_i\big)^{-3/4}$ (\lbl{wind:coefficient}),
      ${\approx}0.0093$ at their fiducial values, with $C_w>1$ meaning wind-dominated. Same
      ordering as Lancaster's $\zeta$ but different exponents; computing both from TRINITY
      output would be a cheap, high-value figure.
\end{itemize}

Their cavity-corrected balance $\tfrac{4\pi}{3}n_i^2(r_i^3-r_w^3)\alpha_B=Q_H$
(\lbl{wind:photoequilibrium}) is one of the three published sources behind
Section~\ref{sec:volume}.

%------------------------------------------------------------------------------
\subsection{Haid et al.\ 2018 --- the sanity anchor. There is no tension.}
\label{sec:haid}

Haid's headline is that the \up{ambient medium} decides the winner, and that the dependence
on it is stronger than on the source: radiation dominates by ${\sim}50\times$ in CNM
($n_0=100$), winds by $10^2$--$10^4$ in WIM ($n_0=0.1$, $T=10^4$\,K), with the switch near
$n_0\sim1\,$cm$^{-3}$. Their $\bar T_{\mathrm{HII}}$ = 7160--8150\,K spans \emph{both}
stellar mass and ambient phase, and they use
$\alpha_B = 2.56\times10^{-13}(T/10^4)^{-0.83}$.

They also measure non-additivity directly:

\begin{litquote}
\small ``In the CNM, the feedback from both processes $p_{\mathrm{Combi}}$ is larger than
$p_{\mathrm{IRad}}+p_{\mathrm{Wind}}$ by $\sim$ 1, 3, and 23 percent\ldots In the WIM, the
difference is a factor of $\sim$ 3.2, 2.8, and
1.9.''\hfill{\scriptsize Haid et al.\ 2018}
\end{litquote}

\begin{keybox}
\textbf{\sffamily Revision 1 claimed this contradicts Lancaster. It does not --- I had
misread \PI.}
\lbl{eq:force\_low\_eta} gives, in the \emph{low}-$\zeta$ limit,
$\dfrac{4\pi\rho_i\Ri^2}{4\pi\rhobar\RSt^2}\approx1+\tfrac12\zeta^3$, and the
surrounding text is unambiguous: \emph{``at small values of $\Req/\RSt$ the momentum from
the CEM solution is actually \textbf{larger} than the momentum given by the separate
idealized solutions.''} The mechanism they give --- \emph{``the wind bubble provides a
volume of gas that is already ionized, liberating LyC photons to ionize gas further out''}
--- is a super-additive coupling, the same direction Haid measure. And \PII\
\lbl{subsubsec:ideal\_sim\_review} places Haid's CNM runs at $\Req/\RSt = 0.35$ \emph{``(all other
values are smaller)''} --- squarely on that branch. (\lbl{app:sim\_comp} is where the input
parameters and the assumed $\ap = 10,\,3,\,1$ for the 12 / 23 / 60 $M_\odot$ cases live.)
\end{keybox}

The ${\sim}35\%$ over-estimate is stated for a \emph{different variable}: \PI\
\lbl{subsec:md\_jfb} gives it at $\Ri/\Rch\approx1$, and since $\Rch=\Req^4/\RSt^3$ that is not the
same condition as $\zeta\approx1$. But do not read that as ``over-estimation belongs only to the
$\Rch$ axis'': immediately after \lbl{eq:force\_low\_eta} \PI\ also says \emph{``it is clear from
[\lbl{fig:numerical\_solution}] that the opposite is true at near-unity values of $\Req/\RSt$''}
--- i.e.\ the sum over-predicts at $\zeta\approx1$ too. The sign change is between \emph{low} $\zeta$
and \emph{near-unity} $\zeta$. \PI's global statement is \emph{``the momentum evolution of the joint
feedback bubble is still within 25\% of the naive value given by the sum of the idealized solutions
in all models presented.''}

\begin{keybox}
\textbf{\sffamily Net, and this is the usable version.} The naive sum is a \up{lower} bound
at low $\zeta$ (Haid's regime, and \PI\ agrees) and an \up{upper} bound near
$\Ri\approx\Rch$ (by ${\sim}35\%$), crossing between. That sign change with regime is
precisely why a coupled closure beats choosing \code{max} or \code{+} globally ---
Section~\ref{sec:closure}.
\end{keybox}

\emph{Also withdrawn:} revision 1 said \PII\ \lbl{subsec:PIR\_cooling\_effect}
``independently confirms Haid's direction''. It does not measure $p_{\mathrm{Combi}}$ vs
$p_{\mathrm{IRad}}+p_{\mathrm{Wind}}$ at all; it measures $\ap$ rising when LyC is on.
Related in spirit, different quantity.

%==============================================================================
\section{What is unambiguously right}

\begin{enumerate}
\item \textbf{Shell momentum equation.} $\mathrm{d}(Mv)/\mathrm{d}t$ form with the
      $\dot M v$ term. Correct, and better than Geen 2019's instantaneous pressure-balance
      closure.
\item \textbf{$R_1$ from free-wind ram balance.} The classical Weaver (1977) inner-shock condition
      $\rho v^2 = \Pb$. \emph{Flagged as outside the verified corpus:} Weaver 1977 is not one of the
      five papers read, and none of them states it --- the only in-corpus form is Lancaster's exact
      strong-shock $P_{\mathrm{hot}}=3\pdotw/(16\pi\Rf^2)$, which is $3/4$ of TRINITY's at the same
      radius. Both conventions are in use; just don't mix them when quoting $\ap$ ---
      Section~\ref{sec:lanc2} gives the exact conversion.
\item \textbf{The transparent cavity.} \code{phi0 = 1} is the standard picture, licensed
      verbatim by Geen 2019 \lbl{winds\_in\_uv}.
\item \textbf{Continuity of the energy$\to$momentum handover.} Because
      $\Pb\to\pdotw/(4\pi R_1^2)$ and $R_1\to R_2$ as $E_b\to0$, $\Pb\to\Pram$
      automatically, so \code{max(P\_thermal, P\_ram)} in the transition phase is a genuine
      continuous handover, not a patch. Your seam ratios 0.995--0.999 confirm it
      \up{for the transition$\to$momentum seam specifically}; \code{PLAN.md} \S3c.1 records
      the others as 0.89--0.92, 0.53--0.96 and 0.86--0.99 for the energy$\to$implicit,
      implicit$\to$transition and regime-switch seams, so quote 0.995--0.999 with its
      own seam attached.
\item \textbf{Self-consistent radiation trapping.} You compute $f_{\mathrm{esc,ion}}$ from
      the actual shell structure instead of Lancaster's threshold-density formula
      (\lbl{eq:phot\_trap\_rhobar}) or Geen 2022's prescribed $r_w(t)$. This directly
      addresses \PI\ \lbl{app:trapping}, which criticises trapping as an \emph{assumption}.
\item \textbf{The confined branch is exactly self-consistent} --- your own \S6b establishes
      this and I agree: hot transparent cavity, skin in pressure equilibrium at $\Pb$, drive
      $=\Pb$ transmitted, no cavity gas mass claimed. Every book balances. It is the
      \emph{driving} branch that does not.
\item \textbf{$\Pext$ and $P_{\mathrm{ISM}}$.} Lancaster \PI\ explicitly drops external
      pressure (\lbl{sec:theory\_joint}: \emph{``we ignore \ldots any external pressures or
      the effects of gravity''}). You keep both, with the right sign. (\PII\ does carry a
      background; the omission is \PI's.)
\item \textbf{Energy bookkeeping of the photoionised term.} The energy ODE charges
      \code{press\_bubble$\cdot$dV} only (\code{energy\_phase\_ODEs.py:274}), so photoionised
      work is not drawn from $E_b$. Correct --- its source is the radiation field,
      continuously resupplied.
\end{enumerate}

%==============================================================================
\section{What can be better --- ranked by \emph{measured} impact}

\subsection{Balance recombination over the ionised layer, not the wind cavity}
\label{sec:volume}

\code{get\_phii\_c3c} sets $n_{\mathrm{C3a}}$ from
$Q_{i,\mathrm{abs}} = \tfrac{4\pi}{3}R_2^3\,\chi_e\alpha_Bn^2$ --- a Str\"omgren balance over
the \up{cavity}. Three papers and TRINITY's own shell solver put the photoionised gas between $r_w$
and $r_i$ (Geen 2019 and Geen 2022 share an author, and Geen 2022's uniform form is the analytic
reduction of its own numerical \lbl{eqn:photoionisation\_equilibrium} --- so: three sources plus your
code, not four independent ones):

\begin{table}[htbp]\centering\small
\begin{tabular}{@{}>{\raggedright\arraybackslash}p{0.42\textwidth}
                  >{\raggedright\arraybackslash}p{0.50\textwidth}@{}}
\toprule
\textbf{\sffamily Source} & \textbf{\sffamily Expression}\\
\midrule
Lancaster \PI\ \lbl{eq:ionreceq2} & $\tfrac{4\pi}{3}(\Ri^3-\Rw^3)\alpha_Bn_{\mathrm{Hi}}^2=Q_0$\\
Geen 2019 \lbl{wind:photoequilibrium} & $\tfrac{4\pi}{3}n_i^2(r_i^3-r_w^3)\alpha_B=Q_H$\\
Geen 2022 \lbl{eqn:photoionisation\_equilibrium\_uniform} & $Q_H=\tfrac{4\pi}{3}n_i^2(r_i^3-r_w^3)\alpha_B$\\
\up{TRINITY} \codet{shell\_structure.py:243} & \codet{\_vol\_ion = R\_IF**3 - rShell0**3}\\
\bottomrule
\end{tabular}
\end{table}

C3a balances over $\tfrac{4\pi}{3}R_2^3$, which in TRINITY's geometry is
$\tfrac{4\pi}{3}r_w^3$ --- \up{exactly the volume all four exclude}. Lancaster add the
physical reason: \emph{``the WBB enhances $\rho_i$ relative to the classical solution, due to
the presence of $\Rw$ in the denominator.''}

The reason to put this first is that it is the \emph{same} defect your Batch 11 measures as
seam C. A cavity Str\"omgren-filled at $n_{\mathrm{C3a}}$ implies
$M_{\mathrm{cav}}/M_{\mathrm{shell}}$ = 0.0952 $\to$ \up{0.5638} by $t=1.5$ on
\code{B3M} --- 57{,}397 vs 101{,}805 $M_\odot$, reproduced by three independent routes ---
while \code{shell\_mass} already equals \up{100.0000\%} of the gas the run has and the winds
inject 54.8 $M_\odot$. On the published picture that gas should not exist at all.

\medskip
Two things I would carry from your own measurements into any write-up. First, B12 showed the
0.5638 is regime-scoped: $M_{\mathrm{cav}}\propto R_2^{3/2}\sqrt{Q_if_{\mathrm{abs}}}$, so
the seam tracks \up{bubble size}, not the degree of HII dominance, and $L_w\times0.1$ gives
0.1296. Always quote it with the config. Second, B11.A's degeneracy result is the strongest
argument for a coupled closure rather than a patch: the photon-conserving fixed point
$x = f_{\mathrm{abs}}(Q_i(1-x))$ has the unique root $x=1$ on 33/33 driving rows, i.e.\ C3a's
own scheme cannot be made photon-conserving without a second equation. Both Geen 2019
\lbl{wind:windpressurebalance} and Geen 2022 \lbl{eqn:PwPibalance} supply exactly that second
equation, and it is the one your shell solver already uses at
\code{shell\_structure.py:125-126}. That is your K5/K6 pair; I would take K6.

%------------------------------------------------------------------------------
\subsection{Shell inertia should exclude the ionised interior --- and this is first-order
\emph{now}}
\label{sec:inertia}

TRINITY puts the entire swept mass inside $R_2$ into the shell. Lancaster \PI\
\lbl{eq:pr\_spitzer\_adj} subtracts the ionised gas:

\begin{litquote}
\small ``A simple adjustment to this momentum inference is to use the shell mass
$\Msh = 4\pi\Ri^3(\rhobar-\rho_i)/3$, that is, to subtract out the mass in ionized gas.
\textbf{Though this reduction is not consistent with the derivation of}
\lbl{eq:HIImomentum2}, we will see that it can be more
accurate.''\hfill{\scriptsize\PI, \lbl{subsec:spitzer}}
\end{litquote}

\PII\ validates it empirically (\lbl{app:spitzer\_momentum},
\lbl{fig:sptiz\_momentum\_comp}) and drops the caveat. \up{Carry the caveat.} The adjustment
is not a consistent re-derivation --- the solution $\mathcal{R}_{\mathrm{Sp}}(t)$ is not
re-derived under the adjusted mass --- so implement it as a debit on the inertia with the
inconsistency stated, exactly as \PI\ does.

\begin{keybox}
\textbf{\sffamily Revision 1 called this ``currently near-irrelevant (your ionised layer is
a thin skin)''. That is false, and your own B11 measured it false.}
On \code{B3M}'s momentum rows $dR_{\mathrm{ion}}/R_2$ = 0.6579--1.3076 with
$dR_{\mathrm{ion}}/dR_{\mathrm{full}}$ median \up{0.9954} --- the momentum shell is
essentially \emph{entirely} ionised. B11.C2 measured the cost: $R_2(t{=}1.5)$ moves
\up{$+8.55\%$} (inertia only) or \up{$+9.22\%$} (inertia and gravity), against an offline
control reproducing the run to $0.871\%$. At $L_w\times0.1$ it falls to $+0.45\%$ / $+0.97\%$,
tracking $M_{\mathrm{cav}}$.
\end{keybox}

So this belongs high in the ordering, not ``the moment the three-radius model lands''.

%------------------------------------------------------------------------------
\subsection{The composition rule: \texttt{max} / \texttt{+} vs a single coupled closure}
\label{sec:closure}

\textbf{Your open question}, from \code{get\_phii\_c3c}'s docstring:
\emph{``$\PCa\propto R_2^{-3/2}$ vs $\Pram\propto R_2^{-2}$: does a real momentum-phase
cavity stay Str\"omgren-filled?''}

Lancaster answer it \emph{for the momentum-driven case}. In the co-evolution phase the
ionised gas is in pressure balance with the wind bubble, so the pressure at the shell is
always the wind-bubble pressure $\ap\pdotw/(4\pi\Rw^2)$ --- photoionisation never adds an
independent pressure. What it changes is \up{where} that pressure acts: \lbl{eq:ionreceq2}
plus \lbl{eq:RiRw\_rel} give $\Ri=\Rw(1+\Rw/\Rch)^{1/3}$, hence (from
\lbl{eq:HIImomentum\_joint1})
\[
\boxed{\;F_b=4\pi\Ri^2 P_i=\ap\pdotw\left(\frac{\Ri}{\Rw}\right)^{\!2}
=\ap\pdotw\left(1+\frac{\Rw}{\Rch}\right)^{\!2/3}.\;}
\]

\begin{keybox}
\textbf{\sffamily Scope, which revision 1 got wrong.} This closure is derived in
\lbl{subsec:md\_jfb} for the \up{momentum-driven phase only}. The energy-driven variant
(\lbl{subsec:ed\_jfb}) replaces $P_i=P_{\mathrm{hot}}$ and requires a separate energy
equation --- four equations in $\Rw$, $\Ri$, $E_w$, $P_{\mathrm{hot}}$. So the numbers below
compare \emph{composition rules against the MD-phase CEM}; they are \emph{not} a
phase-by-phase error budget for TRINITY, and labelling a $-33\%$ column ``1a/1b'' (as
revision 1 did) is a category error.
\end{keybox}

Worth stating plainly too: \PI's pressure balance is \up{imposed, not derived} ---
\emph{``We now imagine that the WBB has come into equilibrium with the PIR with no net force
being applied across its outer edges''}, and \PII: \emph{``the WBB and the PIR are
\textbf{forced} to be at the same pressure.''} It is enforced instantaneously at
$t_{\mathrm{eq}}$, which requires one state variable ($\Ri$) to jump.

I verified numerically that the MD closure has the right limits: it tends to $\ap\pdotw$
(your \emph{confined} branch) as $\Rw/\Rch\to0$ --- 1.0007 at $10^{-3}$ --- and to $\Fsp$
(your \emph{unconfined} branch) as $\Rw/\Rch\to\infty$ --- 1.00007 at $10^4$. So your two
branches \emph{are} the correct asymptotes of the MD CEM. \PI\ is careful that the
strong-wind limit \lbl{eq:approxF2} is stated for $\Rw\ll\Rch$ and the weak-wind limit
\lbl{eq:approxF1} for $\Rw\gg\Rch$; do not claim exactness outside those.

\begin{table}[htbp]\centering\small
\begin{tabular}{@{}cccc@{}}
\toprule
$\Ri/\Rch$ & $F_{\mathrm{sum}}/\Fcem$ & $F_{\max}/\Fcem$\\
\midrule
0.5   & 1.337 & 0.783\\
\up{1.0} & \up{1.342} & \up{0.671}\\
2.0   & 1.318 & 0.772\\
10    & 1.209 & 0.918\\
100   & 1.083 & 0.984\\
\bottomrule
\end{tabular}
\caption{Composition rules against the momentum-driven CEM, at TRINITY's convention
$R_2\equiv\Ri$. The 1.342 at $\Ri=\Rch$ is an independent reproduction of Lancaster's stated
${\sim}35\%$ (\lbl{fig:force\_comp}), which is the check that the normalisation is right.}
\end{table}

\begin{keybox}
\textbf{\sffamily What this licenses you to say.} In the \up{momentum} phase, where TRINITY
uses the bare sum and Lancaster's MD closure applies, $\PHII+\Pram$ over-estimates the
coupled force by up to ${\sim}34\%$, worst at $\Ri\approx\Rch$ --- exactly where your C3c
switch sits. Your own measurement of the same defect from the other side is that the
\code{+P\_ram} term is a $14.1\%$ double-count on \code{B3M} momentum rows
(\code{P\_drive/P\_ram} median $7.095 \to 6.095$ under transmit). The \code{max} column
applies to \emph{energy}-driven phases where the MD CEM does \emph{not} apply, so treat it
as an indication that the two rules bias in opposite directions, not as a quantified error.
\end{keybox}

Also worth knowing, in the weak-wind limit (\lbl{eq:approxF1}):
$F_b\approx\Fsp+\tfrac{\ap\pdotw}{2}\frac{\Rw}{\Ri}$ --- the wind contributes \emph{half} its
momentum flux, further reduced by $\Rw/\Ri$. Your momentum phase credits it the full
$\pdotw$ at the shell radius.

\subsubsection*{Implementation sketch}

You already have the pattern --- \code{solve\_R1} brackets a scalar root on $[0,R_2]$:

\begin{lstlisting}
# given R2 (= R_i, the shell), solve R2 = Rw*(1 + Rw/Rch)**(1/3) for Rw in (0, R2]
# Rch = alpha_B/(12*pi*(mu_H*m_H*ci**2)**2) * (alpha_p*pdot)**2 / Q_ion
#   NOTE: Rch propto alpha_p**2, and Q_ion here must be the ionising rate the
#   *cavity* actually receives (Qi*f_abs), not the raw SPS Q0.
Rw = brentq(lambda x: x*(1.0 + x/Rch)**(1.0/3.0) - R2, 0.0, R2)
P_drive = alpha_p * pdot / (4*np.pi*Rw**2)
\end{lstlisting}

This is a \emph{sketch}, not a drop-in. It is \up{not} backward-compatible, and the edit surface is
larger than the Table-1 row count suggests: there are eight $\Pdrive$ expression sites, of which only
\up{three are live} (\code{energy\_phase\_ODEs.py:253,256} and \code{run\_momentum\_phase.py:445});
the rest are reporting paths that must be kept in step or the force-budget output silently diverges
from the integrated dynamics. In the energy phase you would also keep $\Pb$ from the bubble solve in
place of $\ap\pdotw/(4\pi\Rw^2)$ --- the ED analogue, which needs the energy equation \PI\
\lbl{subsec:ed\_jfb} specifies. Treat it as the shape of K6, and gate it
as you gated C3c.

%------------------------------------------------------------------------------
\subsection{$\ap$: emit it as a diagnostic; do \emph{not} add it as a knob}
\label{sec:alphap}

\begin{keybox}
\textbf{\sffamily Revision 1's \S4.1 was wrong and is withdrawn in full.} It claimed
$\ap\equiv1$ is ``almost certainly why your momentum phase is universally HII-dominated''.
Four things kill that.
\end{keybox}

\begin{enumerate}
\item \textbf{$\ap\equiv1$ in the momentum phase is not a hard-coded omission --- it is what the
      phase asserts.} The force is $4\pi R_2^2 \Pram = \pdotw$ identically, from \code{pRam} plus
      $v_{\mathrm{mech}}=2L/\pdotw_{\mathrm{tot}}$; \code{run\_momentum\_phase.py:587-588}'s
      $R_1=R_2$ (``no inner shock in momentum phase'') states the same assumption for the output
      column, though it is bookkeeping rather than dynamics (Section~\ref{sec:what}, Identity II).
      Lancaster \lbl{eq:alphap\_shock} says $\ap\to1\Leftrightarrow\Rf\to\Rw$, ``the entire bubble is
      made up of the free wind region'' --- the same statement.
\item \textbf{Half of the 4.57--6.82 enhancement is a 3D geometry effect that 1D cannot produce; the
      other half TRINITY already has.} \lbl{eq:alphap\_derive} factorises into a dissipation term
      $\tfrac3{16}\Vw/\langle v_{\mathrm{out}}\rangle$ and a geometry term $4\pi\Rw^2/A_w$
      (Section~\ref{sec:lanc2}). Only the geometry term is identically 1 in 1D. So the honest claim
      is \emph{not} ``a 1D code should have $\ap=1$'' --- TRINITY plainly has $\ap=(R_2/R_1)^2\gg1$
      in its own energy phase --- but: \up{the geometric half is structurally unavailable, and the
      dissipative half is already modelled}, by the bubble energy budget. \PI\
      \lbl{subsubsec:ec\_conditions} says $\ap\approx1$ is expected in the momentum-driven limit
      \emph{``if the correct dissipative scales were resolved''} and \emph{``(at least when magnetic
      fields are not dynamically important)''} --- conditions on resolution and physics, not on
      dimensionality. Any $\ap>1$ imposed by hand would be a calibration of unresolved 3D interface
      geometry, and $A_w$ is exactly what \code{cooling\_boost\_fA} parameterises on the energy side,
      so the two must not be set independently.
\item \textbf{Batch 12 bounds it out.} At $L_w\times0.1$ the momentum phase is 100\%
      HII-dominated at $\PHII/\Pb$ = \up{13.667--14.369}. Inverting that needs
      $\ap\gtrsim14$, twice the largest value \PII\ measures. And at $L_w\times0.1$,
      $\ap\Pram$ is small \emph{whatever} $\ap$ is. So $\ap$ cannot be the explanation.
\item \textbf{It would not even help much at nominal wind.} Because \code{params['Pb']}
      \emph{is} $\Pram$ in phase 2, applying $\ap$ raises $\Pconf$ too. With the measured
      $\PHII/\Pb$ = 5.091--7.156 (median 6.165) and $\ap=6.2$, roughly half the rows flip to
      the confined branch, where $\Pdrive=\ap\Pram=6.2\Pram$ against the shipped
      $7.095\Pram$ --- a ${\sim}13\%$ reduction, not a factor, and a hard branch flip rather
      than a blend. It would also rescale \code{nShell0} $\propto$ \code{params['Pb']}
      (\code{shell\_structure.py:125-126}) and propagate through the whole shell solve. B11.B did not
      measure that scenario, but it measured its close analogue --- replaying each driving row with
      the inner pressure set to $\PCa$ instead of \code{params['Pb']}, a $\times4.70$ (transition) /
      $\times6.17$ (momentum) change in \code{shell\_n0}, comparable in size to $\ap\approx6$ --- and
      found the ionised layer thinning 79--83\% and the dust-absorbed fraction moving
      $0.620\to0.455$ (transition) and $0.607\to0.395$ (momentum). It is not a one-line change.
\end{enumerate}

\begin{figure}[htbp]
\centering
\includegraphics[width=\textwidth]{alphap_origin.pdf}
\caption{\textbf{Left:} TRINITY already computes the quantity comparable to \PII's tabulated $\ap$. Lancaster's
\lbl{eq:alphap\_shock} is $\ap=\tfrac14[3x^2+x^{-2}]$ with $x=\Rw/\Rf$ and
$P_{\mathrm{hot}}=3\pdotw/(16\pi\Rf^2)$; TRINITY's \code{get\_r1} uses the free-wind ram value
$\pdotw/(4\pi R_1^2)$, so at equal bubble pressure $\Rf=(\sqrt3/2)R_1$ and the $4/3$ mismatch
cancels the $3/4$ in the leading term. The solid curve is the exact force ratio
$F/\pdotw=(R_2/R_1)^2$ --- the quantity to emit, and equal to Lancaster's own exact
$4\pi\Rw^2P_{\mathrm{hot}}/\pdotw$; the dashed curve is the full \lbl{eq:alphap\_shock} inversion,
which adds the $x^{-2}$ term Lancaster themselves drop under \emph{``the assumption $\ap\gtrsim1$''}
and which would report $1.1875$ where the force ratio is exactly 1. The momentum phase asserts
$\Rf=\Rw$, pinning it at the left-hand end. \textbf{Right:} what the composition rule costs against
the momentum-driven CEM. The bare sum used in phase 2 over-estimates by up to $34\%$, peaking at
$\Ri/\Rch=0.79$ and $1.342$ at the C3c switch point $R_2=\Rch$; the $\max$ rule under-estimates
there by a comparable amount. Both curves are the MD closure only.}
\label{fig:alphap}
\end{figure}

\subsubsection*{What is actually worth doing}

TRINITY already computes the quantity; it just never names it. Emit, per snapshot, the
\up{force ratio}
\[
\frac{F}{\pdotw} \;=\; \left(\frac{R_2}{R_1}\right)^{2}
\]
through phases 1a, 1b and 1c --- flagged \emph{invalid} inside the \code{dt\_switchon} window, where
the ODE's effective pressure is not $\pdotw/(4\pi R_1^2)$. This is simultaneously TRINITY's exact
force ratio and, under the pressure match, Lancaster's exact $4\pi\Rw^2P_{\mathrm{hot}}/\pdotw$, and
it is what \PII's tabulated $\ap$ approximates (Section~\ref{sec:lanc2}). \up{Do not} emit the exact
\lbl{eq:alphap\_shock} inversion as the headline number: it is correct as algebra but reports
$1.1875$ where the force ratio is exactly 1, which is precisely the regime the 1c$\to$2 handover
approaches. In phase 2, report $1$ by construction and say why.

\begin{keybox}
That gives you a real, publishable comparison against \PII\ \lbl{tab:cem\_comp}.
\textbf{\sffamily The question it answers is not ``is $\ap$ missing'' but ``does TRINITY's
$\ap$ collapse from $(R_2/R_1)^2\gg1$ to 1 too abruptly at the 1c$\to$2 handover?''} --- a
handover governed by \code{ENERGY\_FLOOR = 1e3} in \code{run\_transition\_phase.py:97}, not
by any physical $\ap$ criterion. If $(R_2/R_1)^2$ is still ${\sim}5$ when $E_b$ crosses the
floor, that is a genuine discontinuity in a measurable quantity and a much better paper
result than a fitted knob. If it has already fallen to ${\sim}1$, the momentum phase is
behaving exactly as Lancaster's theory says it should, and the momentum-phase question is
entirely about $\PHII$ (Section~\ref{sec:volume}).
\end{keybox}

Note also: \code{phaseSwitch\_LlossLgain = 0.05} $\Rightarrow$
$\theta = L_{\mathrm{loss}}/L_{\mathrm{gain}} = 0.95$ is the \code{cooling\_balance} trigger
that ends phase 1a (\code{run\_energy\_phase.py:293} reads the threshold, \code{:295} applies it)
and drives the implicit$\to$transition switch (\code{run\_energy\_implicit\_phase.py:1252} reads,
\code{:1298} applies). It is \up{not} the energy$\to$momentum handover, which is
\code{if Eb < ENERGY\_FLOOR} at \code{run\_transition\_phase.py:768}. Revision 1 built an $\ap$ estimator on \lbl{eq:tofap} at $\theta=0.95$; that route
is demoted to a curiosity, and if you use it at all, carry \PI's own conditions --- it
\emph{``assume[s] the modified momentum-driven solution''}, ignores turbulent motions in the
WBB shell, and requires $\ap \gtrsim 1.25$.

%------------------------------------------------------------------------------
\subsection{$T_i$, $\ci$ and $\alpha_B$ are metallicity-blind}
\label{sec:metal}

\code{TShell\_ion = 1e4} and \code{caseB\_alpha = 2.59e-13} are \code{run\_const=True} and
are not derived from \code{ZCloud} (\code{registry.py:396,400} --- no validator, no resolver;
nothing in the tree derives either from metallicity). But:

\begin{itemize}
\item \textbf{Geen 2019} compute $\ci$ from Cloudy as a function of $T_*$, $Z$, $n_i$
      and $\mathcal{U}$
      (\lbl{fig:ionised\allowbreak soundspeed}), at the fiducial values
      $n_i=50$\,cm$^{-3}$ and $\mathcal{U}=-2$. Their text quotes \emph{``on the order of 10 km/s''} --- for gas at
      $10^4$\,K, not explicitly tied to solar $Z$; the higher low-$Z$ value (${\sim}15$ km/s) is
      readable from that figure and a commented-out table, not from quotable body text --- cite it
      as a figure read. Their
      coefficients scale as $C_w\propto \ci^{-3}$ and $C_B\propto \ci^{-4}$
      (\lbl{wind:coefficient}, \lbl{breakout:condition}).
\item \textbf{Geen 2022} carry star-dependent $T_i$ (\lbl{table:app\_starprops}) and
      $\sigma_d = 10^{-21}\,\mathrm{cm}^{-2}\,Z/Z_\odot$.
\item \textbf{Lancaster's} $\zeta\propto \ci^{-1}$ (MD), $\ci^{-3/2}$ (ED).
\item \textbf{Haid} measure $\bar T_{\mathrm{HII}}$ = 7160--8150\,K and use
      $\alpha_B\propto T^{-0.83}$.
\end{itemize}

A factor 1.5 in $\ci$ is a factor ${\sim}3$ in $C_w$ and ${\sim}5$ in $C_B$. Since
\code{ZCloud} already switches SPS tracks and cooling tables, letting it also set $T_i$ (a
small table, or the Geen 2019 fit) closes a real inconsistency --- and moves the answer in
the direction all three papers agree on: \up{low $Z$ makes photoionisation relatively
stronger}. Probably your cheapest genuinely-new physics.

%------------------------------------------------------------------------------
\subsection{Make the overflow transition an explicit event}
\label{sec:overflow}

$Q_{i,\mathrm{abs}}=Q_i f_{\mathrm{abs,ion}}$, so $\PCa\propto\sqrt{f_{\mathrm{abs}}}$ decays
to zero as the shell becomes transparent, and the \code{n\_IF\_Str > 0} gate fails outright at
$f_{\mathrm{abs}}=0$. Meanwhile $\Pext$ is gated on \code{shell\_fAbsorbedIon < 1.0}
(\code{energy\_phase\_ODEs.py:235-238}, and identically in all four runners). So the two respond to
the same physical transition at different points, and over the sequence photoionisation goes from
accelerating to decelerating.

\begin{keybox}
\textbf{\sffamily A correction to revision 1, which had the $\Pext$ gate backwards.} It said the
condition is ``essentially always true, so $\Pext$ is effectively always on''. The opposite holds on
the branch this section is about: \code{shell\_structure.py:230} returns
\code{f\_esc\_ion = max(0.0, ...)}, so $f_{\mathrm{abs}}$ is \emph{exactly} 1.0 whenever the shell
traps everything, and B11.0 measures exactly that on 100\% of transition and 76.5\% of momentum
driving rows (and all of them at $L_w\times0.1$). There $\Pext$ is exactly \up{zero} --- the inward
photoionised term is off precisely where the outward one dominates. (The ISM term is added separately
and unconditionally beyond $r_{\mathrm{cloud}}$, so $\Pext$ is not identically zero out there.)
$f_{\mathrm{abs}}<1$ is instead common in the \emph{energy} phase --- \code{PLAN.md} records
\code{frac\_fabs\_ge\_099} falling to 0.000 on the weak-wind configs --- so the ``always on'' reading
is right only where this section does not apply.
\end{keybox}

Geen 2022 are careful here and should be quoted accurately: \emph{``Reaching the overflow
radius does not guarantee an immediate and strong champagne flow''}; they describe a
rarefaction wave moving back from $r_i$ toward $r_w$ that eventually disperses the shell, and
state \emph{``the precise behaviour of the H\,\textsc{ii} region after it reaches the overflow radius is
mostly beyond the scope of this paper.''} They do \up{not} claim 1D models break there ---
they decline to model it. (Revision 1 said TRINITY's degradation happens ``precisely where
Geen 2022 say a champagne flow begins'', which contradicts the quote it then gave.)

My suggestion is not to model the champagne flow either, but to \up{make the regime change
explicit}: an event in \code{phase\_events.py} on $f_{\mathrm{esc,ion}}$ crossing a threshold,
logged as a distinct fate / \code{SimulationEndCode}, so a run coasting through overflow is
visible in the output instead of looking like ordinary deceleration. Geen 2022's
\lbl{eqn:overflowcondition} gives an analytic overflow radius for $\rho\propto r^{-\omega}$ as
an independent cross-check --- and Section~\ref{sec:geen22}'s reference-mass difference makes
that comparison worth doing rather than assuming agreement.

Related: $\Pext$ is evaluated from the \up{unshocked} profile at $r_{\mathrm{shell}}$
(\code{get\_press\_ion(rShell, ...)}) but applied over $4\pi R_2^2$. Fine for a thin shell;
see Section~\ref{sec:validity}.

%------------------------------------------------------------------------------
\subsection{State the thin-shell validity limit explicitly}
\label{sec:validity}

Your B11.D wording is right and I would put its substance in the method paper verbatim: on
\code{B3M}'s momentum rows $dR_{\mathrm{full}}/R_2$ = 0.6723--1.3078 and
$dR_{\mathrm{ion}}/R_2$ = 0.6579--1.3076, the shell is 99.54\% ionised, and both the ODE's
thin-shell premise and C3a's sharp cavity/shell split are outside their validity range. B12
measured it \up{worse} at low wind ($dR_{\mathrm{ion}}/R_2$ = 1.171--1.438). The trajectory
crosses $dR/R_2\gtrsim1/3$ inside the \emph{transition} phase.

This is not a defect with a proposed fix, and saying so plainly is stronger than not saying
it. It also bounds Section~\ref{sec:overflow}: the $\Pext$ area mismatch is a real error once
$dR/R_2\sim1$.

%------------------------------------------------------------------------------
\subsection{Diagnostics}
\label{sec:diag}

\begin{enumerate}
\item \textbf{Emit $(R_2/R_1)^2$} as in Section~\ref{sec:alphap}, flagged invalid inside the
      \code{dt\_switchon} window and reported as exactly 1 in phase 2, where the code asserts
      $\Rf=\Rw$.
\item \textbf{Emit $\zeta=\Req/\RSt$ and $R_2/\Rch$}, stating the $\ap$ used (both scale with
      it: $\zeta\propto\ap^{1/2}$, $\Rch\propto\ap^2$). These place every TRINITY run
      directly on \PI\ \lbl{fig:feedback\_ratio} / \PII\ \lbl{fig:feedback\_ratio\_comp}.
      Given that C3c's momentum-phase switch \emph{is} $R_2=\Rch$, \code{R2/Rch} is arguably
      the most informative single diagnostic your code could emit.
\item \textbf{\code{F\_ram} is mislabelled in three places.} \code{energy\_phase\_ODEs.py:415},
      \code{run\_energy\_implicit\_phase.py:539} and \code{run\_transition\_phase.py:338} all
      set \code{F\_ram = Pb * 4$\pi$R2**2} --- the \emph{bubble} force, not the ram force ---
      while \code{P\_ram} is reported separately (and is 0 in the energy phase).
      \code{run\_momentum\_phase.py:272} is correct (\code{F\_ram = P\_ram * 4$\pi$R2**2}).
      Reporting only, but it is what every force-budget figure consumes.
\item \textbf{The \code{n\_IF\_Str > 0} gate is vestigial.} Since C3c recomputes the density
      from $R_2$, \code{n\_IF\_Str} and its \code{min(n\_IF\_Str, shell\_n0)} cap no longer
      feed $\PHII$ --- they only gate it. Retire the gate or replace it with the condition it
      intends ($Q_{i,\mathrm{abs}}>0$, $R_2>0$), which \code{get\_phii\_c3c} already checks
      internally. (Same finding as your B11.E.)
\item \textbf{Emit Geen 2019's $C_w$} alongside $\zeta$ --- two independent
      wind/photoionisation orderings from quantities you already carry.
\end{enumerate}

%==============================================================================
\section{On the thin-shell / 1D approximation specifically}
\label{sec:thin}

You flagged this as a constraint. Three observations, one of which reverses revision 1.

\medskip
\textbf{It is not the limitation you might think for $\PHII$.} Everything in
Sections~\ref{sec:volume}--\ref{sec:closure} is achievable in 1D --- Lancaster's entire CEM
\emph{is} a 1D model. Separating $\Rw$ from $\Ri$ costs one algebraic root-find, not a
dimension. Geen 2019/2022 likewise close the whole wind + PIR system in 1D.

\medskip
\textbf{It is the limitation for \emph{half} of $\ap$.} In 1D $A_w\equiv4\pi\Rw^2$, so the
fractal-area half of Lancaster's \lbl{eq:alphap\_derive} is structurally unavailable. The dissipative
half is not: TRINITY carries it as the bubble energy budget, which is why $(R_2/R_1)^2\gg1$ in the
energy phase. So the correct statement is that TRINITY \emph{can} have $\ap>1$ and does, and what it
cannot do is generate the geometric enhancement --- for which any hand-set value would be a
calibration. Your \code{cooling\_boost\_fA} / \code{cooling\_boost\_kappa} knobs are the legitimate
hooks: $f_A$ \emph{is} the $A_w$ factor, on the energy side. So if you ever do calibrate the
geometry, calibrate it once, through $f_A$, and let $\ap$ follow --- do not set the two
independently.

\begin{keybox}
\textbf{\sffamily The thin-shell geometry biases you toward ``confined'' --- but revision 1's
proposed test had the wrong sign.}
Putting \emph{all} swept mass in a thin shell maximises the shell density and hence radiation
trapping, relative to a turbulent, clumpy medium where photons leak through low-density
channels. The porosity form of this criticism is \up{\PII's}
(\lbl{subsec:role\_of\_turbulence}, \lbl{app:supp\_analysis}), which also quantifies the clumping of
the background medium ($\mathfrak{C}=2.39$--10.3, \lbl{subapp:clumping}, measured after the turbulent
evolution but before feedback). \PI\ \lbl{app:trapping} makes a related but distinct argument ---
from background \emph{density gradients} and cluster-vs-single-star geometry --- and itself assumes a
constant-density shell.

Revision 1 suggested a clumping factor $\mathfrak{C}$ multiplying $\alpha_Bn^2$.
\up{That moves the wrong way:} raising the effective recombination rate \emph{shrinks} $\RSt$,
\emph{lowers} $\PCa\propto\RSt^{3/2}$, and makes the confined branch fire \emph{more}, not
less. The correct 1D proxy for photon leakage is a \up{covering / absorbed-fraction} knob ---
an $f_{\mathrm{ion}}<1$ applied to $Q_i$ before the shell solve, representing the fraction of
solid angle over which photons escape through channels. That is also closer to what \PI\
\lbl{app:trapping} actually complains about, and it composes with Geen 2022's $\Omega<4\pi$
geometry (Section~\ref{sec:geen22}).
\end{keybox}

%==============================================================================
\section{Verification}
\label{sec:verify}

Revision 2 went through two adversarial passes after drafting --- one against the papers, one against
the code and the workstream data. Section~\ref{sec:errata} records what revision 1 got wrong;
Section~\ref{sec:pass2} records what those two passes then found in revision 2's own first draft.

\subsection{What was checked, and how}

\begin{table}[htbp]\centering\small
\begin{tabular}{@{}>{\raggedright\arraybackslash}p{0.58\textwidth}
                  >{\raggedright\arraybackslash}p{0.36\textwidth}@{}}
\toprule
\textbf{\sffamily Claim} & \textbf{\sffamily Result}\\
\midrule
\codet{solve\_R1} root $\Rightarrow \Pb=(\gamma-1)\tfrac{3}{4\pi}L/(vR_1^2)$;
  $=\pdotw/(4\pi R_1^2)$ at $\gamma=5/3$ & exact (SymPy \codet{simplify == 0})\\
\codet{pRam} with $v_{\mathrm{mech}}=2L/\pdotw_{\mathrm{tot}}$
  $\Rightarrow F_{\mathrm{ram}}=\pdotw_{\mathrm{tot}}$ & exact\\
C3c crossover (momentum phase) $=\Req^4/\RSt^3=\Rch$ & ratio $=1$ exactly\\
$\ap$ from Weaver-convention $R_2/R_1$ via \lbl{eq:alphap\_shock} &
  $=(R_2/R_1)^2+\tfrac{3}{16}(R_1/R_2)^2$ exactly; $+1.16\%$ at $R_2/R_1{=}2$, $+0.23\%$ at 3\\
$\ap=1\Leftrightarrow R_2/R_1=\sqrt3/2$ in TRINITY's convention & $0.8660254$\\
C3c crossover vs $f_{\mathrm{abs}}$ & $R_2^{\mathrm{cross}}\propto f_{\mathrm{abs}}^{-1}$
  (SymPy solve), i.e.\ \up{outward}\\
$\Fcem\to\ap\pdotw$ as $\Rw/\Rch\to0$ & 1.00067 at $10^{-3}$\\
$\Fcem\to\Fsp$ as $\Rw/\Rch\to\infty$ & 1.000067 at $10^{4}$\\
$F_{\mathrm{sum}}/\Fcem$ at $\Ri=\Rch$ & \up{1.342} (Lancaster state ${\sim}1.35$)\\
$F_{\max}/\Fcem$ at $\Ri=\Rch$ & 0.671\\
\lbl{eq:tofap} inverted for $\ap$ & round-trips $\theta$ to $10^{-6}$\\
$\zeta(\ap=6.25)$, \PII\ cluster & 0.991 (paper quotes ${\approx}0.98$)\\
Lancaster's exact force $\tfrac34x^2$ under the pressure match & $=(R_2/R_1)^2$ exactly\\
C3c crossover $=\Rch$ requires $\ap=1$ & crossover $=\Rch/\ap^2$ in general\\
$F_{\mathrm{sum}}/\Fcem$ true maximum & 1.344 at $\Ri/\Rch=0.79$, not at 1.00\\
\bottomrule
\end{tabular}
\end{table}

Code claims re-checked against the staged source this revision:

\begin{itemize}\small\raggedright
\item \code{run\_momentum\_phase.py} --- \code{:587-588} and \code{:893} ($R_1=R_2$);
      \code{:272} (\code{F\_ram} from \code{P\_ram});
      \code{:585}, \code{:669}, \code{:891} (\code{Pb = pRam}).
\item \code{energy\_phase\_ODEs.py:415}, \code{run\_energy\_implicit\_phase.py:539},
      \code{run\_transition\_phase.py:338} --- all three set \code{F\_ram} from \code{Pb}.
\item \code{get\_bubbleParams.py:365} (the \code{params['Pb']} comparison) and
      \code{:495-503} (the \code{dt\_switchon} ramp).
\item \code{shell\_structure.py} --- lines \code{119-120}, \code{125-126}, \code{158},
      \code{182}, \code{243}, \code{253}.
\item \code{run\_energy\_phase.py:55} (\code{TFINAL\_ENERGY\_PHASE = 3e-3}) and \code{:293};
      \code{run\_energy\_implicit\_phase.py:1252};
      \code{run\_transition\_phase.py:97} (\code{ENERGY\_FLOOR = 1e3});
      \code{registry.py:396,400,407}.
\end{itemize}

Every \LaTeX{} label was re-checked against the cited paper, including which of
\lbl{eq:Rch\_def} / \lbl{eq:Rch\_Req\_rel} / \lbl{eq:Rch\_app} carries which statement.

%------------------------------------------------------------------------------
\subsection{What revision 1 got wrong}
\label{sec:errata}

Recorded rather than quietly patched, because the point of the exercise was that each claim
be true.

\begin{table}[htbp]\centering\footnotesize
\begin{tabular}{@{}c>{\raggedright\arraybackslash}p{0.42\textwidth}
                  >{\raggedright\arraybackslash}p{0.46\textwidth}@{}}
\toprule
\textbf{\sffamily \#} & \textbf{\sffamily Revision 1 claim} & \textbf{\sffamily Status}\\
\midrule
1 & ``$\ap\equiv1$ is almost certainly why your momentum phase is HII-dominated''; add an
    \codet{alpha\_p} knob & \up{WRONG.} Phase 2 applies $F=\pdotw$ identically via \codet{pRam};
    Batch 12 bounds the required factor at $\gtrsim14$\\
2 & ``Do \emph{not} report $(R_2/R_1)^2$; it converges to $\tfrac43\times$ Lancaster's
    $\ap$'' & \up{WRONG, backwards.} With Weaver's convention the $4/3$ cancels:
    $\ap=(R_2/R_1)^2+\tfrac3{16}(R_1/R_2)^2$\\
3 & Haid's non-additivity ``goes the opposite way to Lancaster's analytic result'' &
    \up{WRONG.} \lbl{eq:force\_low\_eta} gives super-additivity at low $\zeta$; \PII\ puts
    Haid at $\zeta\le0.35$ --- same branch\\
4 & ``Geen \& de Koter 2022 --- you already implement their model'' & \up{WRONG chronology.}
    Geen 2022 cites Rahner 2017. Common Draine (2011) inheritance\\
5 & ``You are strictly ahead'' of Geen 2022 & \up{WRONG.} They carry $Z$-dependent
    $\sigma_d$ and star-dependent $T_i$\\
6 & Shell-mass adjustment is ``currently near-irrelevant (thin skin)'' & \up{WRONG, measured
    false.} $dR_{\mathrm{ion}}/dR_{\mathrm{full}}$ median 0.9954; worth $+8.55\%/+9.22\%$\\
7 & The $-33\%$ / $+34\%$ table labelled ``1a/1b'' and ``momentum'' & \up{WRONG scope.} The
    closure is MD-phase only\\
8 & Partial absorption shifts the crossover by $f_{\mathrm{abs}}^{1/3}$ & \up{WRONG} in
    exponent and sign: $f_{\mathrm{abs}}^{-1}$, outward\\
9 & ``\PI\ levels three criticisms at \textsc{warpfield}'' & \up{WRONG.} Eight; the omitted
    ones include its central thesis\\
10 & ``\codet{frac HII-dominated = 0.0000} on every energy and implicit row of all five
    configs'' & \up{STALE.} Batch 7 records \codet{B3MW001} at 78.4\% HII-dominated in the
    energy phase\\
11 & ``\codet{F\_ram} is mislabelled in two phases'' & \up{WRONG count.} Three: energy,
    implicit, transition\\
12 & Seam ratios ``0.995--0.999'' quoted for the handover generally &
    \dn{UNDER-QUALIFIED.} That is the transition$\to$momentum seam; others are 0.89--0.92,
    0.53--0.96, 0.86--0.99\\
13 & Identity I stated without qualification & \dn{INCOMPLETE.} Holds for stored
    \codet{Pb} always; for the ODE-effective pressure only outside the
    \codet{dt\_switchon} window\\
14 & Clumping factor $\mathfrak{C}$ on $\alpha_Bn^2$ & \up{WRONG SIGN.} It makes the confined
    branch fire \emph{more}. Use $f_{\mathrm{ion}}<1$ on $Q_i$\\
15 & ``$\theta=0.95$ is your transition trigger'' used as an energy$\to$momentum handover
    value & \dn{MISLABELLED.} It ends 1a and drives 1b$\to$1c; the handover is
    \codet{Eb < ENERGY\_FLOOR = 1e3}\\
16 & ``the CEM ignores \ldots on-the-fly bubble cooling \ldots TRINITY has all seven'' &
    \dn{OVERSTATED.} The ED-CEM does solve an energy equation; six, and what it lacks is a
    dissipation model\\
17 & \PII\ ``independently confirms Haid's direction'' & \dn{OVERSTATED.} Different quantity
    ($\ap$, not $p_{\mathrm{Combi}}$ vs the sum)\\
18 & Geen's overflow criterion identified with TRINITY's & \dn{IMPRECISE.} TRINITY matches
    Geen's \emph{analytic} form ($M(<r_w)$); it overflows earlier by $(r_i/r_w)^{3-\omega}$\\
19 & ``$\PCa/\Pram=3.8$--7.6'' & \dn{SUPERSEDED} by measurement: 5.091--7.156, median 6.165
    (\codet{B3M} momentum); 13.667--14.369 at $L_w\times0.1$\\
20 & \PI's CEM pressure balance described as a result & \dn{IMPRECISE.} Imposed:
    \emph{``we now imagine''}, \emph{``forced to be at the same pressure''}\\
21 & Table cited \codet{energy\_phase\_ODEs.py:385} as the 1c site & \dn{HALF RIGHT.} \codet{:385}
    really is unreachable, but \codet{run\_transition\_phase.py:331} is \emph{reporting only} --- the
    live 1c site is \codet{energy\_phase\_ODEs.py:253}\\
22 & ``three different thresholds'' at overflow & \dn{TWO} --- but see Section~\ref{sec:pass2}: the
    $\Pext$ half was itself wrong\\
\bottomrule
\end{tabular}
\end{table}

%------------------------------------------------------------------------------
\subsection{What the \emph{second} verification pass changed}
\label{sec:pass2}

Revision 2 was itself run past two independent adversarial passes before release. Nine more findings,
all incorporated above; the four that changed a conclusion:

\begin{table}[htbp]\centering\footnotesize
\begin{tabular}{@{}>{\raggedright\arraybackslash}p{0.36\textwidth}
                  >{\raggedright\arraybackslash}p{0.58\textwidth}@{}}
\toprule
\textbf{\sffamily Draft-2 claim} & \textbf{\sffamily Status}\\
\midrule
``$\Pext$ is effectively always on'' & \up{REFUTED.} $f_{\mathrm{abs}}$ saturates at exactly 1.0 on
  100\% of transition and 76.5\% of momentum driving rows, so $\Pext$ is exactly \emph{zero} there.
  The gate is on in the \emph{energy} phase instead\\
``$\ap\equiv1$ is forced by \codet{R1 = R2}'' & \up{REFUTED as causal.} \codet{params['R1']} is read
  in exactly one place (\codet{get\_bubbleParams.py:115}, energy/implicit code); in phase 2 it is
  bookkeeping. What forces $F=\pdotw$ is \codet{pRam} with $v_{\mathrm{mech}}=2L/\pdotw$\\
``a 1D code cannot produce $\ap>1$ by this mechanism at all'' & \up{REFUTED.}
  \lbl{eq:alphap\_derive} has \emph{two} factors; only the geometric one is 1 in 1D. TRINITY has
  $\ap=(R_2/R_1)^2\gg1$ in its own energy phase\\
``\codet{run\_transition\_phase.py:331} is the live 1c site'' & \up{REFUTED.} The transition RHS
  delegates to \codet{get\_ODE\_Edot\_pure}; the live sites are \codet{energy\_phase\_ODEs.py:253,256}
  and \codet{run\_momentum\_phase.py:445}\\
\bottomrule
\end{tabular}
\end{table}

Five smaller ones: the $\Req$ definition is \lbl{eq:RE\_MD\_def}, not \lbl{eq:eta\_def}; the
$\zeta=0.35$ statement is in \lbl{subsubsec:ideal\_sim\_review}, not \lbl{app:sim\_comp};
\lbl{eq:force\_low\_eta} has no $\ci^2$; the ramp discrepancy is $3.31\times$, not $3.2\times$; and
\PI's eight ``distinct'' criticisms are better grouped as six, with two I had missed
(constant-density PIR, and shell inertia).

Two numbers I could not source and therefore removed: the lower bound \code{0.487} on Batch 7's
confinement ratio (it appears only in revision 1 of this document, and
\code{data/b7\_confinement\_screen.csv} is not in my staged copy).

\subsection{What I did \emph{not} verify}

\begin{itemize}
\item \textbf{Anything requiring a TRINITY run.} Every measured number in this revision is
      taken from your committed \code{data/*.csv} and the Batch 11/12 write-ups, cited as
      such. I re-derived none of them.
\item \textbf{Whether $(R_2/R_1)^2$ actually reaches \PII's 4.6--6.8 before the
      1c$\to$2 handover.} That is the central open question of Section~\ref{sec:alphap} and
      it is unmeasured. The cheapest test is an offline pass over existing snapshots --- no
      solver run.
\item \textbf{Geen 2019's low-$Z$ $\ci\approx15\,$km/s.} Readable from
      \lbl{fig:ionisedsoundspeed} and a commented-out table; the body text quotes only
      \emph{``on the order of 10 km/s''}.
\item \textbf{\PII's $\ap$ at higher resolution.} The MWR sequence
      $5.57 \to 6.20 \to 6.82$ is not converged and the paper says so.
\end{itemize}

\textbf{Staleness caveat.} \code{docs/dev/} carries its own ``may be out of date'' banner, and
its \code{LITERATURE\_ASSESSMENT.md} is revision 1 of \emph{this} document, under a maintainer
note that a correction was inbound. This is that correction; its \S4 items should be re-read
against Section~\ref{sec:errata} before any of them is acted on. Every code claim here was
re-checked against source at the 2026-08-18 pull; \code{trinity/} was byte-identical to the
copy read for revision 1.

%==============================================================================
\section{Suggested order of work}

\begin{enumerate}
\item \textbf{Emit the diagnostics} (Section~\ref{sec:diag}):
      $\ap=(R_2/R_1)^2+\tfrac3{16}(R_1/R_2)^2$ through 1a/1b/1c, $\zeta$, $R_2/\Rch$, Geen's
      $C_w$. Zero risk, no solver run, and it answers ``where do my runs sit on
      \lbl{fig:feedback\_ratio}''.
\item \textbf{Offline: does $(R_2/R_1)^2$ reach 4.6--6.8 before the 1c$\to$2 handover?}
      (Section~\ref{sec:alphap}) One pass over existing snapshots. This is the question
      revision 1 should have asked.
\item \textbf{Run B11.G} --- score the shipped closure against Geen 2019
      \lbl{wind:photoequilibrium} + \lbl{wind:windpressurebalance} / Geen 2022
      \lbl{eqn:photoionisation\_equilibrium\_uniform} + \lbl{eqn:PwPibalance} on your own
      trajectory. Cheap, offline, and the direct input to Section~\ref{sec:volume}.
\item \textbf{The balance volume} (Section~\ref{sec:volume}) --- K5 as the minimal move, K6
      as the right one. Pre-register the ablation as you did for C3c; it will move fates.
\item \textbf{Shell-mass adjustment} (Section~\ref{sec:inertia}) --- already measured at
      $+8.6$--$9.2\%$; land it with Lancaster's own consistency caveat stated.
\item \textbf{$Z$-dependent $T_i$, $\ci$, $\alpha_B$} (Section~\ref{sec:metal}).
\item \textbf{Paper-scale:} three-radius model, the coupled closure across all four phases,
      an $f_{\mathrm{ion}}$ leakage knob (Section~\ref{sec:thin}), and the stated thin-shell
      validity limit (Section~\ref{sec:validity}) in the methods section.
\end{enumerate}

%==============================================================================
\section*{Sources}
\addcontentsline{toc}{section}{Sources}

\subsubsection*{Papers (project documents)}
\begin{itemize}\small\raggedright
\item \code{lancaster2025.tex} --- Lancaster, J.-G.\ Kim, Bryan, Menon, Ostriker \& C.-G.\ Kim,
      \emph{The Co-Evolution of Stellar Wind-blown Bubbles and Photoionized Gas I}.
\item \code{lancaster2025b.tex} --- \emph{\dots II: 3D RMHD Simulations and Tests of
      Semi-Analytic Models}.
\item \code{geen2019.tex} --- Geen, Pellegrini, Bieri \& Klessen, \emph{When H\,\textsc{ii}
      Regions are Complicated}.
\item \code{geen2022.tex} --- Geen \& de Koter, \emph{Bottling the Champagne}.
\item \code{Haid118\_arxiv.tex} --- Haid, Walch, Seifried, W\"unsch, Dinnbier \& Naab,
      \emph{The relative impact of photoionizing radiation and stellar winds on different
      environments}.
\end{itemize}

\subsubsection*{TRINITY source (read-only)}
\begin{itemize}\small\raggedright
\item \code{trinity/bubble\_structure/get\_bubbleParams.py}, \code{bubble\_luminosity.py}
\item \code{trinity/phase1\_energy/energy\_phase\_ODEs.py}, \code{run\_energy\_phase.py}
\item \code{trinity/phase1b\_energy\_implicit/run\_energy\_implicit\_phase.py}
\item \code{trinity/phase1c\_transition/run\_transition\_phase.py}
\item \code{trinity/phase2\_momentum/run\_momentum\_phase.py}
\item \code{trinity/shell\_structure/shell\_structure.py}, \code{get\_shellODE.py}
\item \code{trinity/sps/update\_feedback.py}, \code{trinity/\_input/registry.py}
\end{itemize}

\subsubsection*{Workstream documents}
\begin{itemize}\small\raggedright
\item \code{docs/dev/phii-identity/PLAN.md} --- \S0 (C-0 carve-out), \S3c.1, \S6b, Batch 7,
      Batch 11 (B11.0, B11.A--D), Batch 12, \S7.1 candidate register (K1--K9).
\item \code{docs/dev/phii-identity/README.md} --- \S5, \S6.
\item \code{docs/dev/phii-identity/LITERATURE\_ASSESSMENT.md} --- revision 1 of this document,
      with the B11 cross-check.
\item \code{data/b7\_confinement\_screen.csv}, \code{b10\_wind\_profile.csv},
      \code{b11\_mass\_ledger.csv}, \code{b11\_photon\_ledger.csv},
      \code{b11\_mass\_dynamics.csv}, \code{b12\_lowwind\_mass\_ledger.csv},
      \code{b12\_lowwind\_mass\_dynamics.csv}.
\end{itemize}

\end{document}
